% !TEX root = main.tex
\section{Background}
\label{sec:background}

Short-term retail demand forecasts support replenishment, inventory allocation,
markdowns, and labour planning. The data typically consist of many related
item--location series, often with intermittent sales and information about
prices, promotions, and the calendar. This combination has encouraged the use
of global models, which learn across series, as well as a long tradition of
simple per-series benchmarks.

\subsection{Forecasting methods}
\label{sec:taxonomy}

We distinguish four broad model families. Statistical methods such as
exponential smoothing, ARIMA, Theta, Seasonal Naive, Croston, and TSB are fitted
to individual series \citep{croston1972forecasting,teunter2011intermittent}.
They remain useful benchmarks: a more complex model should demonstrate an
improvement over a forecast that repeats the most recent seasonal pattern.

Gradient-boosted trees such as LightGBM \citep{ke2017lightgbm} and XGBoost
\citep{chen2016xgboost} are commonly fitted globally across items and stores.
Lagged demand, calendar indicators, prices, promotions, and item or store
identifiers allow the model to combine temporal and cross-sectional
information. This family dominated the M5 competition and is an important
practical comparator for retail applications. Multi-step forecasts can be
produced recursively with one model or directly with a separate model for each
horizon. As shown later, this choice can materially affect the result.

Neural forecasting methods also learn across series. They include recurrent
models such as DeepAR \citep{salinas2020deepar}, attention-based models such as
TFT and PatchTST \citep{lim2021tft,nie2023patchtst}, and architectures such as
N-BEATS and N-HiTS \citep{oreshkin2020nbeats,challu2023nhits}. They require
training on the target collection and therefore differ from the zero-shot
models at the centre of this paper.

Foundation models are pre-trained on large collections of time series and can
be transferred to a new dataset without estimating new model parameters.
Table~\ref{tab:fm-scope} lists the models covered by the literature extraction
or the experiments. Our empirical comparison is restricted to zero-shot use;
fine-tuned results are recorded in the data but excluded from the main
contrasts.

\begin{table}[htbp]
\centering
\caption{Foundation models covered by the study.}
\label{tab:fm-scope}\label{tab:fm_models}
\footnotesize
\begin{tabularx}{\textwidth}{l l X X X}
\toprule
Model & Parameters & Architecture & Version considered & Reference \\
\midrule
Chronos family & 9--710\,M & Chronos / Chronos-2 & Chronos-2 120M; Bolt Tiny/Base &
  \citet{ansari2024chronos,ansari2025chronos2} \\
TimesFM~2.5 & 200\,M & Decoder-only & 2.5 checkpoint & \citet{das2025timesfm2} \\
Moirai~2.0 & 11--305\,M & Quantile model & Small checkpoint & \citet{liu2025moirai2} \\
TiRex & \textasciitilde{}35\,M & xLSTM & Public checkpoint & \citet{auer2025tirex} \\
TabPFN-TS & \textasciitilde{}12\,M & Tabular PFN & fev-bench tasks & \citet{mueller2025tabpfnts} \\
Lag-Llama & 6\,M & Llama backbone & Public checkpoint & \citet{rasul2023laglama} \\
TimeGPT & Not disclosed & Proprietary & Published API results & \citet{garza2023timegpt} \\
\bottomrule
\end{tabularx}
\end{table}

\subsection{Retail forecasting benchmarks}
\label{sec:existing-benchmarks}

Three benchmark sources are particularly relevant. The M5 competition
\citep{makridakis2022m5} contains 30,490 daily Walmart item--store series and
covariates for calendar events, prices, and the US Supplemental Nutrition
Assistance Program. Its official measure, \WRMSSE{}, scales errors by each
series' historical variability and weights economically important series more
heavily. Gradient-boosted trees were used by the leading entries.
\label{sec:m5-benchmark}

M5 also illustrates why apparently minor evaluation choices matter. A large
share of item--day observations are zero. Mean per-series \WAPE{} gives the
same weight to a low-volume item as to a high-volume item and becomes unstable
when test-period demand is close to zero. Aggregate \WAPE{} and \WRMSSE{}
place different weights on the series and may therefore produce different
rankings. We keep these measures separate throughout the analysis.

GIFT-Eval \citep{aksu2024gift} was designed to reduce the risk that a public
test set had appeared in a foundation model's pre-training data. It covers 23
datasets from several domains and reports skill relative to Seasonal Naive
with bootstrap intervals. Only part of the suite is retail-specific, so we
extract results at the task level rather than use its overall ranking.

Fev-bench \citep{shchur2025fevbench} contains 100 forecasting tasks, including
46 with exogenous variables. It covers foundation models, statistical
methods, and tree baselines, and is consequently useful for studying whether
benchmark conclusions change with the comparator or accuracy measure. We use
the 20 tasks identified as retail-related.

Forecasting competitions provide strong comparisons because models share data
and metrics, but each competition represents a limited set of demand patterns.
Reviews of forecasting and deep learning provide broader context but either
precede the recent foundation-model literature or do not compare these models
with retail tree baselines on common panels \citep{lim2021tft}. The evidence
assembled here connects those two forms of evaluation: broad public benchmark
coverage and narrower experiments in which the forecasts can be paired at the
series level.
